\section{The varying case when \texorpdfstring{$b_2=4$}{b2=4}}
\label{var:section}
\label{sec:variation}

\begin{proposition}[The varying rank-four case]\label{var:exclusion}
Let $X$ be a projective irreducible holomorphic symplectic manifold of
dimension $2n$ with $b_2(X)=4$, and let
\[
                         f:X\longrightarrow B,\qquad
                         g:X\longrightarrow C
\]
be projective Lagrangian fibrations with normal projective bases.
Assume that the polarized family of smooth fibers of $g$ has nonzero
variation and that, for every smooth fiber $A_t$ of $g$, the restriction
$f|_{A_t}:A_t\to B$ is finite and surjective.
Fix a connected minimum stratum $S_0\subset B$ and its smooth
Deligne--Mumford stack $\cS$ as in Propositions~\ref{str:minimum}
and~\ref{str:products}.  Then
\[
 H^{a,b}(\cS)=0\quad\text{for }a\ne b,\qquad
 \chi_{\mathrm{top}}(S_0)=\sum_p h^{p,p}(\cS)>0.
\]
\end{proposition}

Proposition~\ref{secsys:rank-four-variation} supplies these two
fibrations when the original rank-four fibration has nonzero polarized
variation, with the first fibration still mapping to the original
base $B$.  The positive Euler characteristic in the proposition
contradicts Proposition~\ref{eu:euler-zero} when the fixed stratum has
a singular maximal chart.

Fix the data in the proposition, and let $C^\circ$ be the connected
smooth open of regular values of $g$.
Resolve the reduced inverse image of $S_0$ in $g^{-1}(C^\circ)$,
retaining the components which dominate $C^\circ$, and restrict to a
connected dense algebraic open $T\subset C^\circ$ over which the
resolution is smooth.  The base $T$ is smooth and quasi-projective.
We obtain smooth projective families
\[
 \pi:\mathcal A=g^{-1}(T)\longrightarrow T,
 \qquad \rho:\mathcal Q\longrightarrow T,
\]
with maps $j:\mathcal Q\to\mathcal A$ and $h:\mathcal Q\to\cS$.
The latter map uses the normalized maximal-chart covers, which are
\'etale over the smooth source by purity; it therefore retains all the
stabilizer data of $\cS$.  Put $d=\dim S_0$.  Each $Q_t$ has dimension
$d$, possibly with several components, and its map to $S_0$ is dominant and generically
finite.  Indeed $A_t\to B$ is finite.  The proper closed images of the
omitted nondominating components can be removed from $T$.
Put
\[
 V_\Q=R^1\pi_*\Q,
 \qquad U^i_\Q=R^i\rho_*\Q,
 \qquad H_1=V^{1,0}.
\]
An origin on the fibers $A_t$ is unnecessary: their holomorphic
one-forms are translation invariant and trivialize their cotangent
bundles.  All variations in this section are taken over the same
algebraic open $T$.

To prove diagonal cohomology, suppose that
$0\ne\alpha\in H^{a,b}(\cS)$ with $a>b$.  We will turn its
pullbacks into a flat tensor whose exterior factor has a nonzero
all-holomorphic component.  We will find a joint monodromy lowering
operator whose block from $V^{1,0}$ to $V^{0,1}$ is invertible,
whereas its action on the factor $U^b$ vanishes after $b+1$
lowerings.  Invariance of the tensor will
then force its $a$th exterior lowering to vanish, contradicting the
all-holomorphic component.  The finite maps $A_t\to B$ produce the
tensor; the rank-four invariant-cycle count, together with the
nonzero variation of $g$, produces the invertible lowering operator.

\begin{lemma}[Differentials and the top cup component]
\label{var:cup-image}
For every $t\in T$ there is a natural factorization
\begin{equation}
 h_t^*\Omega^1_{\cS}\longrightarrow
 j_t^*\Omega^1_{A_t}\longrightarrow\Omega^1_{Q_t}.
 \label{var:factorization}
\end{equation}
If $0\ne\alpha\in H^{a,b}(\cS)$, its pullback is a nonzero global
flat section $y$ of $U^{a+b}_\C$.  Moreover, for the morphism of
variations
\begin{equation}
 \mathcal T:\Lambda^a V_\Q\otimes U^b_\Q\longrightarrow U^{a+b}_\Q,
 \qquad v\otimes u\longmapsto j^*v\smile u,
 \label{var:cup-map}
\end{equation}
we have, at every point of $T$,
\begin{equation}
 y_t\in\mathcal T_t\bigl(
       \Lambda^a H_{1,t}\otimes (U^b_t)^{0,b}\bigr).
 \label{var:top-cup}
\end{equation}
\end{lemma}

\begin{proof}
We first check the independence needed to define
\eqref{var:factorization}.  On a maximal chart, lift a form on the
smooth orbit to an ordinary ambient differential, using the local
product presentation of Proposition~\ref{str:products}.  Every tangent
vector to the smooth finite cover of this chart maps, at a point above
the orbit, into the orbit tangent space.  Here is the trace argument
in this finite-map situation.  Work with a pointed finite local branch
$R\subset E$ of generic degree $m$, with $E$ regular.  Extend a tangent
vector upstairs to a local derivation $\delta$ of $E$.  The rule
\[
 r\longmapsto m^{-1}\operatorname{Tr}_{\operatorname{Frac}(E)/
                            \operatorname{Frac}(R)}(\delta r)
\]
is a derivation of $R$: its values are regular by normality, and the
Leibniz rule follows from $R$-linearity of trace.  Its value at the
marked point is the differential image of the chosen tangent vector,
since all conjugates in the pointed branch have the same residue and
trace on the residue is multiplication by $m$.  Evaluations of
ambient derivations are precisely the orbit tangent directions, by
Proposition~\ref{str:minimum}.

The normalized pullback of a maximal chart to $A_t$ is smooth and
\'etale over $A_t$.  Indeed the quasi-\'etale exceptional locus has
codimension at least two, as does its inverse image under the finite
map $A_t\to B$, and purity applies.  Thus the preceding tangent
calculation applies there.  The difference of two lifts of an orbit
form vanishes on the orbit tangent, so its pullback vanishes as an
ambient cotangent vector at every point above the orbit.  Restricting
to the reduced inverse image and then pulling to its resolution kills
this difference.  The resulting maps are compatible on the chart
overlaps and with the finite group actions; hence they descend to
\eqref{var:factorization}.

The map of the entire family to the fixed stack $\cS$ makes $h^*\alpha$
a global flat section.  It is nonzero on every fiber.  To see this,
choose $\beta\in H^{d-a,d-b}(\cS)$ with $\alpha\smile\beta\ne0$, using
Poincar\'e duality and the Hodge decomposition on $\cS$.  Pullback on
$H^{2d}(\cS,\C)$ is nonzero: the $d$th power of a coarse ample class
has positive integral on a dominating component of $Q_t$.  This also
holds for $d=0$, where the class is the unit.  Since $\cS$ is connected,
its top cohomology is one-dimensional; consequently
$h_t^*(\alpha\smile\beta)\ne0$ and $h_t^*\alpha\ne0$.

Finally, take the $a$th exterior power of
\eqref{var:factorization}, and use
\[
 j_t^*\Omega^a_{A_t}
       =\Lambda^a H^0(A_t,\Omega^1_{A_t})\otimes\cO_{Q_t}.
\]
On $H^b$ the resulting sheaf maps factor Dolbeault pullback through
\[
 \Lambda^a H^0(A_t,\Omega^1_{A_t})\otimes H^b(Q_t,\cO_{Q_t}).
\]
Under the Dolbeault identification the last arrow is cup product
with the pulled-back holomorphic forms.  This is exactly
\eqref{var:top-cup}.
\end{proof}

Choose flat polarizations on the variations in
\eqref{var:cup-map}, and the induced tensor polarization on its
source.  Write $\mathcal T^\dagger$ for the adjoint.  This adjoint is
flat, although the Hodge metrics themselves generally vary.  Indeed,
both sides of \eqref{var:cup-map} have weight $a+b$, and a morphism of
Hodge structures commutes with their Weil operators.  The adjoint
for their positive Hodge metrics is therefore its transpose for the
flat polarization forms.  It is a morphism of variations.
Consequently
\begin{equation}
 z=\mathcal T^\dagger y
      \in\Gamma(T,\Lambda^a V_\C\otimes U^b_\C)
 \label{var:flat-z}
\end{equation}
is flat.  Its component
\begin{equation}
 z_{\mathrm{top}}\in
      \Lambda^a H_1\otimes(U^b)^{0,b}
 \label{var:z-top}
\end{equation}
is nonzero at every point.  In fact, choose $w$ in this component
with $\mathcal T w=y$.  If the orthogonal projection of $z$ to this
component were zero, then
\[
 0=\langle z,w\rangle
   =\langle y,\mathcal T w\rangle
   =\langle y,y\rangle>0,
\]
a contradiction.  Orthogonality here follows from the factor Hodge
decompositions of the tensor Hodge metric.

\begin{lemma}[The simultaneous Higgs operator]
\label{var:joint-monodromy}
Let $T$ be a connected smooth quasi-projective algebraic variety, and
let $\pi:\mathcal A\to T$ and $\rho:\mathcal Q\to T$ be smooth
projective morphisms.  The fibers of $\rho$ may be disconnected.
For an integer $b\ge0$, put
$V_\Q=R^1\pi_*\Q$ and $U^b_\Q=R^b\rho_*\Q$, choose flat
polarizations on these rational variations, and write $H_1=V^{1,0}$.
Assume the following fixed-part condition for every pure tensor word
in $V_\Q,U^b_\Q$ and their duals: its full-monodromy invariant subspace
is a constant Hodge structure, so every Hodge component of a global
flat invariant tensor is again globally flat.  This condition holds
for the geometric systems just specified, as verified below.

Let $G$ be the full Zariski closure of the monodromy on
$V_\Q\oplus U^b_\Q$.  Then $G$ is linearly reductive, though it may
be disconnected.  Its Lie algebra $\mathfrak g_\C$, transported by
conjugation as a flat subbundle of
$\End(V_\C)\oplus\End(U^b_\C)$, is a subvariation of Hodge structure
of weight zero.  For every tangent vector $x\in T_tT$,
\begin{equation}
     (\theta_{V,x},\theta_{U,x})\in\mathfrak g_t^{-1,1}.
 \label{var:higgs-in-lie}
\end{equation}
The image
\begin{equation}
 E_t=\im\bigl(\mathfrak g_t^{-1,1}
                      \longrightarrow\Hom(H_{1,t},\overline H_{1,t})\bigr)
           \subset\Sym^2 H_{1,t}^*
 \label{var:E}
\end{equation}
is a subbundle preserved by the Chern connection.  In unitary frames it
satisfies
\begin{equation}
       A B^* C+C B^* A\in E_t\qquad(A,B,C\in E_t),
 \label{var:triple}
\end{equation}
where $B^*$ is the Hermitian adjoint matrix.
\end{lemma}

\begin{proof}
First we verify the stated fixed-part condition from the geometric
hypotheses.  Write $U=U^b$.  A pure tensor word has the form
\[
 \mathbb M=
 V^{\otimes r}\otimes(V^\vee)^{\otimes r'}
 \otimes U^{\otimes s}\otimes(U^\vee)^{\otimes s'}.
\]
The polarizations give $V^\vee\simeq V(1)$ and
$U^\vee\simeq U(b)$, so
\[
 \mathbb M\simeq
 \bigl(V^{\otimes(r+r')}\otimes U^{\otimes(s+s')}\bigr)(r'+bs').
\]
The untwisted word is a flat Hodge direct summand of
$R^{r+r'+b(s+s')}\varpi_*\Q$ for the corresponding fiber product
$\varpi:\mathcal Y\to T$ of copies of $\pi$ and $\rho$.  The canonical K\"unneth
inclusion and projection commute with monodromy and Hodge
decomposition; no algebraic K\"unneth correspondence is required.
The product remains smooth and projective when fibers of $\rho$ are
disconnected.  Deligne's Theorem~4.1.1(ii) and Corollary~4.1.2 of
\cite{Deligne1971} supply invariant lifts from a nonsingular
compactification of this product and make their Hodge components
globally flat.  Applying the K\"unneth projection and then the Tate
twist proves the fixed-part condition for $\mathbb M$.
It also proves it for flat Hodge direct summands such as exterior
and symmetric powers.  The summand projectors decompose every tensor
construction on $V\oplus U$ into these pure words.  Since $T$ is
connected, global flat sections are exactly full-monodromy invariants,
and taking the Zariski closure does not change them.

The same geometric hypotheses give semisimplicity of
$V_\Q\oplus U^b_\Q$ by Deligne's Proposition~4.2.5 and
Theorem~4.2.6 of \cite{Deligne1971}.  Both cohomology systems have
rational structures, integral lattices modulo torsion, and
polarizations on the smooth algebraic base $T$.  Semisimplicity and
faithfulness imply that the unipotent radical of $G^\circ$ is trivial.
The component group $G/G^\circ$ is finite, so averaging complements
in characteristic zero makes the full, possibly disconnected group
$G$ linearly reductive.

A linearly reductive subgroup of a general linear group is the pointwise
stabilizer of its invariant mixed tensors.  One way to check the
pointwise assertion is to realize the subgroup as the stabilizer of
a line in a rational representation, realize that representation in
mixed tensors, and choose an invariant complementary subspace to the
line by linear reductivity.  Its projection is an invariant tensor; fixing
that tensor preserves the line and hence belongs to the subgroup.
Applying this to the faithful representation $V\oplus U^b$, with its
invariant summand projectors included, gives
\begin{equation}
 \mathfrak g_t=\{D: D\tau=0\text{ for every invariant mixed tensor }
                                      \tau\text{ at }t\}.
 \label{var:tensor-annihilator}
\end{equation}

For a flat invariant tensor of a single Hodge type, the three parts
of its connection have different Hodge degrees.  Its lowering part
therefore vanishes separately.  Applying this to every flat Hodge
component just described shows that the simultaneous pair in
\eqref{var:higgs-in-lie} annihilates every tensor in
\eqref{var:tensor-annihilator}.  This proves
\eqref{var:higgs-in-lie}; it is an assertion about the Lie algebra of
the pointwise stabilizer.  Also, the invariant tensor spaces are
Hodge graded, so their common infinitesimal annihilator is Hodge
graded.  It is flat by conjugation transport, and therefore is a
subvariation of the indicated weight-zero endomorphism variation.

The grading-preserving part of its connection is the restriction of
the ambient Chern connection.  Projection to $\End(V)$ is a morphism
of variations; its image is a subvariation, and its degree $(-1,1)$
part is $E$.  This proves constant rank and Chern parallelism.  The
weight-one polarization identifies $\overline H_1$ with $H_1^*$ and
makes these lowering matrices symmetric.

For completeness, the bracket calculation takes place in the
polarization-preserving Lie algebra.  In a unitary Hodge frame a
lowering operator has block matrix
\[
 L_A=\begin{pmatrix}0&0\\ A&0\end{pmatrix}
       \quad\text{on }H_1\oplus\overline H_1.
\]
The conjugate of a lift of $L_B$ in $\mathfrak g^{-1,1}$ is in
$\mathfrak g^{1,-1}$ and has upper block $B^*$, up to the same
irrelevant polarization sign for all $B$.  The double bracket of
lifts of $L_A$, this upper operator, and $L_C$ has lower block
$AB^*C+CB^*A$, up to sign.  Its Hodge degree is $(-1,1)$; projection
to the first factor proves \eqref{var:triple}.
\end{proof}

Apply the lemma to the geometric systems constructed above; its
fixed-part condition was verified in the proof.  The Higgs matrices
in $E$ have an additional symmetry supplied by
the Lagrangian system.  Symplectic contraction gives an isomorphism
\begin{equation}
 \xi:T_T\xrightarrow{\ \sim\ }H_1,
       \qquad x\longmapsto[\iota_{\widetilde x}\sigma|_{A_t}],
 \label{var:contraction}
\end{equation}
independent of the lift $\widetilde x$ since the fiber is Lagrangian.
Regard $\xi$ as a $V_\C$-valued one-form.  It satisfies
$d_\nabla\xi=0$.  Indeed its pairing with any locally transported
fiber one-cycle is the fiber integral of the closed form $\sigma$;
exterior differentiation commutes with this integration, and the
restriction of $\sigma$ to the fibers is zero.  Taking the
$V^{0,1}$ component gives
\[
       \theta_{V,x}(\xi(y))=\theta_{V,y}(\xi(x)).
\]
The polarization already gives symmetry in the two matrix slots.
After the identification \eqref{var:contraction}, these two
symmetries give a cubic
\begin{equation}
 c_t\in\Sym^3H_{1,t}^*,\qquad
          c_t(x,-,-)=\theta_{V,\xi^{-1}(x)}\in E_t.
 \label{var:cubic}
\end{equation}
This is the Donagi--Markman cubic associated with a Lagrangian system
\cite[Section~7.2, Lemmas~7.1--7.2 and Remark~7.3]{DonagiMarkman1996};
the calculation above gives the symmetry in the conventions used here.
Nonzero variation means that this cubic is not identically zero.

\begin{lemma}[Maximal-rank kernels and a symmetric cubic]
\label{var:matrix}
Let $H$ be a finite-dimensional Hermitian vector space and let
$E\subset\Sym^2 H^*$ be a complex subspace satisfying
\eqref{var:triple} in unitary coordinates.  Suppose
$c\in\Sym^3H^*$ and $c(x,-,-)\in E$ for every $x\in H$.
If $A\in E$ has the maximum rank occurring in $E$, then
\[
                c(k,-,-)=0\qquad(k\in\ker A).
\]
\end{lemma}

\begin{proof}
Let $r=\rk A$; the assertion is immediate if $r=0$.  By unitary
symmetric diagonalization, take coordinates in which
$A=\operatorname{diag}(\lambda_1,\ldots,\lambda_r,0)$, with
$\lambda_i>0$.  The triple operation gives all odd powers
$\operatorname{diag}(\lambda_i^{2m+1},0)$.  Interpolation on the
finitely many nonzero singular values therefore gives
\[
          P=\operatorname{diag}(I_r,0)\in E,
          \qquad H=F\oplus K,\quad K=\ker A.
\]
The linear map $Z\mapsto PP^*Z+ZP^*P$ preserves $E$ and has
eigenvalues $2,1,0$ on the $FF$, off-diagonal, and $KK$ blocks.
Its polynomial spectral projections thus separate these blocks
inside $E$.  A nonzero $KK$ block would, when added to $P$, have rank
larger than $r$.  Hence every member of $E$ has zero $KK$ block.
Write
\[
 \mathcal D=\left\{D:
       \begin{pmatrix}D&0\\0&0\end{pmatrix}\in E\right\},\qquad
 \mathcal B=\left\{B:
       \begin{pmatrix}0&B\\B^{\mathsf t}&0\end{pmatrix}\in E\right\}.
\]
Here $D$ is symmetric and $B:K\to F$ is a rectangular matrix.
Triple closure, first with $(P,D,P)$ and then with
$(D,P,\left(\begin{smallmatrix}0&B\\B^{\mathsf t}&0\end{smallmatrix}\right))$,
gives
\begin{equation}
                 D^*\in\mathcal D,\qquad DB\in\mathcal B.
 \label{var:block-stability}
\end{equation}

For $B\in\mathcal B$, the Schur complement computes
\[
 \rk\begin{pmatrix}I_r&B\\B^{\mathsf t}&0\end{pmatrix}
                 =r+\rk(B^{\mathsf t}B).
\]
Maximality of $r$ implies $B^{\mathsf t}B=0$.  Polarizing in $B$
shows that, for a fixed $k\in K$, the space
\[
                W_k=\{Bk:B\in\mathcal B\}\subset F
\]
is totally isotropic for the complex symmetric form
$(u,v)\mapsto u^{\mathsf t}v$.  In detail,
$B^{\mathsf t}C+C^{\mathsf t}B=0$, and evaluation on $(k,k)$ gives
$(Bk)^{\mathsf t}(Ck)=0$.  Equation~\eqref{var:block-stability}
shows that $W_k$ is stable under every $D\in\mathcal D$ and every
$D^*$.

Put $S=c(k,-,-)$.  Every cubic contraction has zero $KK$ block;
by symmetry, for $f\in F$ and $k'\in K$,
\[
                 c(k,f,k')=c(f,k,k')=0.
\]
Thus $S$ has only an $FF$ block, also denoted by $S$, in
$\mathcal D$.  For $f\in F$, let $B_f$ be the off-diagonal block of
$c(f,-,-)$.  Cubic symmetry gives $Sf=B_fk$, so
\begin{equation}
                  \im S\subset W_k.
 \label{var:S-image}
\end{equation}
Since $S$ is symmetric and $W_k$ is isotropic, for $w\in W_k$ and
$f\in F$ we have
\[
                  (Sw)^{\mathsf t}f=w^{\mathsf t}Sf=0.
\]
It follows that $S$ annihilates $W_k$.  But $S^*$ preserves $W_k$
by \eqref{var:block-stability}; together with
\eqref{var:S-image} this proves
\begin{equation}
                           SS^*S=0.
 \label{var:positive-triple-zero}
\end{equation}
The singular values of $SS^*S$ are the cubes of those of $S$.
Hence $S=0$.  This is the Hermitian positivity step; no positivity
of the complex symmetric form was asserted.
\end{proof}

\begin{lemma}[Existence of an invertible lowering matrix]
\label{var:invertible}
Under the hypotheses of Proposition~\ref{var:exclusion}, every fiber
of the bundle $E$ attached to the families $\pi,\rho$ above contains
an invertible matrix.
\end{lemma}

\begin{proof}
Let $r$ be the maximum rank in $E_t$.  It is independent of $t$,
because Chern parallel transport is unitary and carries $E_t$ onto
$E_{t'}$ by congruence.  Define intrinsically
\begin{equation}
 K_t^{\mathrm{all}}=
      \sum_{\substack{A\in E_t\\\rk A=r}}\ker A\subset H_{1,t}.
 \label{var:all-kernels}
\end{equation}
Unitary Chern transport preserves the entire collection of these
kernels, and therefore their span.  Thus \eqref{var:all-kernels}
is a global Chern-parallel subbundle of constant rank.  No matrix
$P$, chosen kernel, or finite cover enters its definition.

By Lemma~\ref{var:matrix}, the cubic vanishes when any argument lies
in $K^{\mathrm{all}}$.  In particular all lowering Higgs operators
annihilate this subbundle.  Write the real flat connection on $V$ as
\[
                       \nabla=D+\theta+\overline\theta,
\]
with $D$ the Chern connection.  On $K^{\mathrm{all}}\subset V^{1,0}$,
$\theta$ vanishes by the cubic and $\overline\theta$ vanishes by
Hodge degree.  Complex conjugation in the real variation shows that
$\overline\theta$ vanishes on
$\overline{K^{\mathrm{all}}}\subset V^{0,1}$, while $\theta$ vanishes
there by degree.  Consequently
\[
       K^{\mathrm{all}}\oplus\overline{K^{\mathrm{all}}}
\]
is the complexification of a real flat subbundle $L\subset V_\R$;
its flat connection is unitary.  Positivity of the Hodge metric
shows that the alternating polarization is nondegenerate on $L$.
Its symplectic orthogonal $L^\perp$ is also a global flat subbundle.
These statements concern the full flat connection on $T$, and hence
full monodromy around every loop.

We now compute the relevant full-monodromy invariant space.
The use of the one-dimensional fiber-restriction image to prove
irreducibility over $\R$ also appears in Bakker's proof of Matsushita's
variation conjecture \cite[Lemma~5]{Bakker2025}, building on Voisin's
argument. Here it rules out the specific flat summand extracted from
the maximal-rank kernels.  By the
global invariant-cycle theorem, in the form of
\cite[Theorem~4.1.1(ii)]{Deligne1971} applied to
$g^{-1}(C^\circ)\to C^\circ$ with smooth projective compactification
$X$, the image of
\[
       H^2(X,\Q)\longrightarrow H^2(A_t,\Q)=\Lambda^2 V_{\Q,t}
\]
is precisely the invariant subspace.  Its dimension is one: the
independent classes $\sigma$, $\overline\sigma$, and $g^*c_1(H_C)$
restrict to zero, whereas an ample class restricts nontrivially;
here $H_C$ is ample and $b_2(X)=4$.
The inclusion $T\hookrightarrow C^\circ$ induces a surjection of
fundamental groups, because loops in a smooth complex variety can
be moved off a closed subset of positive complex codimension.
Since $V|_T$ is the restriction of the same local system, its
monodromy image is unchanged.  Therefore
\begin{equation}
            \dim_\R(\Lambda^2 V_{\R,t})^{\pi_1(T,t)}=1.
 \label{var:one-invariant}
\end{equation}
Invariants commute with extension from $\Q$ to $\R$: they are the
common kernels of rational monodromy-minus-identity maps, and a
finite subcollection suffices in finite dimension.

If $L$ were nonzero and proper, the polarization composed with the
two flat projections onto $L$ and $L^\perp$ would give two independent
real invariant alternating forms.  The polarization identifies the
dual alternating-square representation with $\Lambda^2 V_\R$, so
this contradicts \eqref{var:one-invariant}.  Rationality of $L$ is
not needed.  If $L=V_\R$, the Higgs field vanishes identically,
contradicting the chosen nonzero variation of $g$ (which remains
nonzero on every dense open).  Thus $L=0$.  If $r<n$, every
maximum-rank matrix would have a nonzero kernel, making
$K^{\mathrm{all}}$ nonzero.  Hence $r=n$.
\end{proof}

\begin{proof}[Proof of Proposition~\ref{var:exclusion}]
Suppose $0\ne\alpha\in H^{a,b}(\cS)$ with $a>b$.  The flat section
$z$ in \eqref{var:flat-z} is fixed by the joint algebraic monodromy
$G$, and its component \eqref{var:z-top} is nonzero.  At a point
$t\in T$, choose an invertible matrix $A\in E_t$ and a lift
$\nu\in\mathfrak g_t^{-1,1}$.  Let $N_V$ be its induced action on
$\Lambda^a V_t$ and $N_U$ its action on $U^b_t$, acting on their
respective tensor factors.  Invariance gives
\[
                         (N_V+N_U)z_t=0.
\]
Since $U^b_t$ has Hodge types $(p,b-p)$ with $0\le p\le b$, the
operator $N_U$ satisfies $N_U^{b+1}=0$.  The two tensor actions
commute, so induction gives
\begin{equation}
                    N_V^a z_t=(-1)^aN_U^a z_t=0,
 \label{var:nilpotence}
\end{equation}
because $a>b$.

On $V_t$, the operator $\nu$ maps $H_{1,t}$ isomorphically onto
$\overline H_{1,t}$, and kills $\overline H_{1,t}$.  Consequently
$N_V^a$ kills every summand of $\Lambda^a V_t$ having fewer than
$a$ factors in $H_{1,t}$.  On the whole sector with all $a$ exterior
factors holomorphic, including its $U^b_t$ factor, it is
\[
 \left.N_V^a\right|_{\Lambda^a H_{1,t}\otimes U^b_t}
   =a!\,\Lambda^a A\otimes\id_{U^b_t}:
 \Lambda^a H_{1,t}\otimes U^b_t
   \xrightarrow{\ \sim\ }
 \Lambda^a\overline H_{1,t}\otimes U^b_t.
\]
The projection of $z_t$ to this sector is nonzero, because its further
projection to $(U^b_t)^{0,b}$ is the nonzero component
\eqref{var:z-top}.  The displayed map is injective and every other
exterior sector is killed, so $N_V^a z_t\ne0$,
contradicting \eqref{var:nilpotence}.  Thus no $H^{a,b}(\cS)$ with
$a>b$ is nonzero; complex conjugation excludes $a<b$ as well.

It follows that all odd cohomology of $\cS$ vanishes and
\[
 \chi_{\mathrm{top}}(S_0)
     =\sum_p h^{p,p}(\cS)>0,
\]
since $h^{0,0}(\cS)=1$.  Rational cohomology of the smooth stack is
that of its coarse space, including when its stabilizers are
ineffective.  This proves the stated positive Euler characteristic.
\end{proof}
